% ============================================================
% M3 S2 导学件·波1 片件 —— 课时08 直线与圆的位置关系（2.3.3）
% 生成：M3 成卷轮 S2 波1 臂3｜2026-09-14｜编译 cwd＝本目录（中文路径禁 \input 绝对路径）
%   xelatex main-true.tex ×2 → 含详解印本；xelatex main-false.tex ×2 → 纯题版（单源双本）
% 输入链（全只读）：题面库 课时08-直线与圆的位置关系.md（题面逐字装配）＋ -答案侧.md（值逐字回嵌）
%   ＋ 定稿/批B-课时08-直线与圆的位置关系.md（详解回嵌源）；toolchain＝qp-m3.sty
%   （钉后版 md5 7c3930362be8a0a2bdf21bbf8ac16573，本地挂载副本，破损即停）。
% 母版体例：照抄 成卷/导学件/课时01-坐标法/（骨架/宏用法/符号路由 preamble 整块照抄）。
% 片级注记承接（批B 台账下发）：①08-T15 题面【注】＝抽验销项③（导数符号分支需补 a<−1 排除，
%   答案不变；销项落点在批B 定稿 T15 亲算行 0913 已落盘），题面注为元数据不入印面；
%   ②08-T16 源含示意图，印面去图（批B 归属栏同注），全部数学信息在文字中，无信息损失。
% 键账：件内 ansblock 键集 ≡ 题面库片键集 ≡ manifest 键序（21 键，零缺漏零多余）；
%   预习位零键零答案（口令④裁定前口径）；印面号 1—21＝探究 1—16＋课堂评价 17—21，
%   键↔号映射见 值台账-课时08-直线与圆的位置关系.json。
% 括线判据（承重墙禁放宽）：估高 top-5（wlen 全角1·ASCII0.5，剔\命令／23 栏宽字口径）
%   → 括线模，逐键读数入值台账；其余灰底。本片括线键见 值台账 items 判模列。
% 门谱读数：见 过程件 工作区/_tmpM3S2波1臂3_0914/（编译 log、门谱输出、目验 PNG 双档）。
% ============================================================
\documentclass[fontset=none]{ctexart}
\usepackage{qp-m3}
% —— 印面逐字符号路由补钉（照抄课时01 母版 0914 实证块；禁改 sty 本体保 md5） ——
\xeCJKDeclareCharClass{Default}{"2212}%
\setCJKfamilyfont{nscblk}[Path=C:/提示词/工作区/字替对照-0909/variantF/fonts/,BoldFont=NSC-w600.ttf]{NSC-w500.ttf}%
\xeCJKDeclareSubCJKBlock{nscblk}{"25B1}%
\newcommand{\pxparallelogram}{{\CJKfamily{nscblk}▱}}%
% —— 断行微调（照抄母版：CJK 胶收缩 0.01→0.05em；\emergencystretch 不设） ——
\xeCJKsetup{CJKglue={\hskip 0.10em plus 0.08em minus 0.05em}}%
% —— 页脚件名（qp-headfoot 复用入口） ——
\renewcommand{\qpzhangming}{第二章\quad 平面解析几何}%
\renewcommand{\qpjianming}{导学件}%

\begin{document}

\zhangtitle{第二章\quad 平面解析几何}
\jietitle{2.3.3 直线与圆的位置关系}
\mubiaoline
\mubiaomu{1}{掌握直线与圆相交、相切、相离三种位置关系的几何法（\(d\)与\(r\)比较）与代数法（联立方程组）判定，能处理含参、恒成立与集合语言的位置判定问题；}
\mubiaomu{2}{掌握圆的切线方程求法（过圆上点、过圆外点、切两轴定圆）与切线长公式，会处理切线条数计数与切线长最值问题；}
\mubiaomu{3}{掌握弦长公式与圆心角的几何计算，能将视线遮挡、阿波罗尼斯圆等实际问题几何化为直线与圆的位置关系问题.}

\vspace{1.7mm}
\begin{multicols}{2}
\emergencystretch=1em

% ============================================================
% 课前预习（知识点导学填空；零键零答案——预习键位按检查口令④裁定前口径，
% \kongbai 空线＋\zhentib 判断题面形，两档等形不泄答）
% ============================================================
\huaxing{课}{前}{预}{习}{知识导学\quad 素养初识}

\zsd{一}{直线与圆的位置关系}

\tiaomu{1}{直线\(l\)与圆\(C\)的位置关系有相交、相切、相离三种．几何法：圆心\(C\)到\(l\)的距离\(d\)与半径\(r\)比较；代数法：联立直线与圆的方程，由交点个数判定．若\(d<r\)则\kongbai{}，若\(d=r\)则\kongbai{}．\zhuzhu{含参位置判定优先几何法（\(d\)与\(r\)比大小）；恒成立问题把\(d\)表为参数的函数后求最值再比.}}

\zsd{二}{圆的切线}

\tiaomu{1}{过圆\(x^{2}+y^{2}=r^{2}\)上一点\(P(x_{0},y_{0})\)的切线方程为\kongbai{}；过圆外一点作圆的切线有\kongbai{}条，切线长为\(\sqrt{|PC|^{2}-r^{2}}\)．\zhuzhu{过圆外一点求切线防漏斜率不存在的支（先验点在圆外再分类）；切线长最小值＝“点圆距平方减半径平方”开方，点在定直线上时化为圆心到直线的距离.}}

\zsd{三}{弦长与圆心角}

\tiaomu{1}{弦心距\(d\)、半弦长、半径构成直角三角形：弦长\(|AB|=2\sqrt{r^{2}-d^{2}}\)；弦长、半径与圆心角\(\theta\)满足\(|AB|=2r\sin(\theta/2)\)．\zhuzhu{垂径定理是弦长问题的核心：弦心距垂直平分弦；中点弦的方向垂直于圆心与弦中点的连线.}}

\zhenhead{判断正误(正确的打\gou{},错误的打\cha{})}

\zhentib{(1)}{若圆心到直线的距离等于半径，则直线与圆相切．}

\zhentib{(2)}{过圆上一点可作圆的两条切线．}

\zhentib{(3)}{直线与圆相交时，圆心到直线的距离小于半径．}

\liubai[9mm]{此处书写}

% ============================================================
% 课中探究（考点探究〔例+变式〕＝练习槽 16 键；题后紧跟答案＝ansblock 内嵌，
% 值照答案侧逐字，详解承批B 亲算回嵌＋印面清洗；装配序＝印面号 1—16）
% ============================================================
\huaxing{课}{中}{探}{究}{考点探究\quad 素养小结}

% ---------- 探究点一 直线与圆的位置关系 ----------
\tjdnr{一}{直线与圆的位置关系}{例\textbf{1}}{中档(过定点直线)}{直线\(l\)：\(y=k(x-1)+1\)和圆\(x^{2}+y^{2}-2y=0\)的位置关系为\nobreak (\kongwei)}

\bindopt {\fontsize{10.5pt}{\qpTJgLeadLiti}\selectfont \makebox[\dimexpr0.5\linewidth-1em\relax][l]{A．相交}\makebox[\dimexpr0.5\linewidth-1em\relax][l]{B．相切或相交}\par}

\bindopt {\fontsize{10.5pt}{\qpTJgLeadLiti}\selectfont \makebox[\dimexpr0.5\linewidth-1em\relax][l]{C．相离}\makebox[\dimexpr0.5\linewidth-1em\relax][l]{D．相切}\par}

\begin{ansblock}[2章-练-课时08-T5]
% ans:2章-练-课时08-T5
\ansitem{1}{A}
\ansnote{详解}{圆\(x^{2}+(y-1)^{2}=1\)，圆心\((0,1)\)。直线\(y=k(x-1)+1\)过定点\((1,1)\)（\(x=1\)时\(y=1\)），恰为圆心，故\(d=0<r\)，恒相交，故选：A．}
\end{ansblock}

\liB{变式\textbf{2}}{简单(含参恒判定)}{}{直线\(x\cos\theta+y\sin\theta=r\)和\(x^{2}+y^{2}=r^{2}\)的位置关系是\kongbai{}．}
\xiexwei{9mm}

\begin{ansblock}[2章-练-课时08-T11]
% ans:2章-练-课时08-T11
\ansitem{2}{相切}
\ansnote{详解}{圆心\((0,0)\)，半径\(r\)。\(d=|-r|/\sqrt{\cos^{2}\theta+\sin^{2}\theta}=r\)，恒相切．}
\end{ansblock}

\liB{变式\textbf{3}}{中档(含参与坐标轴公共点)}{}{若圆\(C\)：\(x^{2}+y^{2}-2mx+2\sqrt{m}\cdot y+2=0\)与\(x\)轴有公共点，则实数\(m\)的取值范围是\kongbai{}．}
\xiexwei{9mm}

\begin{ansblock}[2章-练-课时08-T12]
% ans:2章-练-课时08-T12
\ansitem{3}{[√2, +∞)}
\ansnote{详解}{表圆：\(\sqrt{m}\)有意义即\(m\geqslant 0\)，\(r^{2}=m^{2}+m-2>0\)，在\(m\geqslant 0\)下\(m>1\)。\par\noindent 与\(x\)轴有公共点\(\iff|\sqrt{m}|\leqslant r\)，即\(m\leqslant m^{2}+m-2\)，即\(m^{2}\geqslant 2\)，得\(m\geqslant\sqrt{2}\)。故\(m\in[\sqrt{2},\ +\infty)\)．}
\end{ansblock}

\liB{变式\textbf{4}}{中档(视线应用·相离条件)}{}{已知圆\(C\)：\((x-1)^{2}+y^{2}=1\)，点\(A(-2,0)\)及点\(B(3,a)\)，从点\(A\)处观察点\(B\)，要使视线不被圆\(C\)挡住，则实数\(a\)的取值范围为\kongbai{}.}
\xiexwei{9mm}

\begin{ansblock}[2章-练-课时08-T13]
% ans:2章-练-课时08-T13
\ansitem{4}{(−∞, −5√2/4)∪(5√2/4, +∞)}
\ansnote{详解}{视线不被挡\(\iff\)直线\(AB\)与圆相离。\(k_{AB}=a/5\)，直线：\(ax-5y+2a=0\)。\par\noindent \(d=|3a|/\sqrt{a^{2}+25}>1\)，即\(9a^{2}>a^{2}+25\)，\(8a^{2}>25\)，\(|a|>5/(2\sqrt{2})=5\sqrt{2}/4\)．}
\end{ansblock}

\tjdnr{一}{直线与圆的位置关系}{例\textbf{5}}{中档(集合语言·交点计数)}{已知集合\(A=\{(x,y)\mid x^{2}+y^{2}=1\}\)，集合\(B=\{(x,y)\mid y=|x|-1\}\)，则集合\(A\cap B\)的真子集的个数为\nobreak (\kongwei)}

\bindopt {\fontsize{10.5pt}{\qpTJgLeadLiti}\selectfont \makebox[\dimexpr0.25\linewidth-0.5em\relax][l]{A．\(3\)}\makebox[\dimexpr0.25\linewidth-0.5em\relax][l]{B．\(4\)}\makebox[\dimexpr0.25\linewidth-0.5em\relax][l]{C．\(7\)}\makebox[\dimexpr0.25\linewidth-0.5em\relax][l]{D．\(8\)}\par}

{\ansblockgrayfalse
\begin{ansblock}[2章-练-课时08-T8]
% ans:2章-练-课时08-T8
\ansitem{5}{C}
\ansnote{详解}{\(x\geqslant 0\)：\(y=x-1\)，\(x^{2}+(x-1)^{2}=1\)，即\(2x^{2}-2x=0\)，\(x=0\)（\(y=-1\)）或\(x=1\)（\(y=0\)）；\(x<0\)：\(y=-x-1\)，\(x^{2}+(x+1)^{2}=1\)，即\(2x^{2}+2x=0\)，\(x=-1\)（\(y=0\)）（\(x=0\)不在此段）。\par\noindent 交点\((0,-1)\)、\((1,0)\)、\((-1,0)\)共3个，真子集\(2^{3}-1=7\)个，故选：C．}
\end{ansblock}}

\xiaojie{①识别：判定直线与圆（含参直线、含参圆、恒成立）的位置关系、集合语言交点计数、实际相离（视线）条件，判归本题型.}
\bindp {\kaishu ②操作：几何法优先——圆心到直线距离\(d=|Ax_{0}+By_{0}+C|/\sqrt{A^{2}+B^{2}}\)与\(r\)比大小；含参先守表圆前提，恒成立转参数最值.}
\bindp {\kaishu ③收束：代数法判别式作复核；交点个数与位置关系互译（相离0、相切1、相交2），实际问题先建模为直线与圆再判定.}

% ---------- 探究点二 圆的弦长问题 ----------
\tjdnr{二}{圆的弦长问题}{例\textbf{6}}{中档(弦长公式)}{直线\(l\)：\(3x+4y-1=0\)被圆\(C\)：\(x^{2}+y^{2}-2x-4y-4=0\)所截得的弦长为\nobreak (\kongwei)}

\bindopt {\fontsize{10.5pt}{\qpTJgLeadLiti}\selectfont \makebox[\dimexpr0.25\linewidth-0.5em\relax][l]{A．\(2\sqrt{5}\)}\makebox[\dimexpr0.25\linewidth-0.5em\relax][l]{B．\(4\)}\makebox[\dimexpr0.25\linewidth-0.5em\relax][l]{C．\(2\sqrt{3}\)}\makebox[\dimexpr0.25\linewidth-0.5em\relax][l]{D．\(2\sqrt{2}\)}\par}

\begin{ansblock}[2章-练-课时08-T1]
% ans:2章-练-课时08-T1
\ansitem{6}{A}
\ansnote{详解}{\((x-1)^{2}+(y-2)^{2}=9\)，圆心\((1,2)\)，\(r=3\)。\(d=|3+8-1|/5=2\)。弦长\(=2\sqrt{9-4}=2\sqrt{5}\)，故选：A．}
\end{ansblock}

\liB{变式\textbf{7}}{简单(圆内点最短弦)}{}{已知圆\(x^{2}+y^{2}=9\)的弦过点\(P(1,2)\)，当弦长最短时，该弦所在直线方程为\nobreak (\kongwei)}

\bindopt {\fontsize{10.5pt}{\qpTJgLeadLiti}\selectfont \makebox[\dimexpr0.5\linewidth-1em\relax][l]{A．\(x+2y-5=0\)}\makebox[\dimexpr0.5\linewidth-1em\relax][l]{B．\(y-2=0\)}\par}

\bindopt {\fontsize{10.5pt}{\qpTJgLeadLiti}\selectfont \makebox[\dimexpr0.5\linewidth-1em\relax][l]{C．\(2x-y=0\)}\makebox[\dimexpr0.5\linewidth-1em\relax][l]{D．\(x-1=0\)}\par}

\begin{ansblock}[2章-练-课时08-T2]
% ans:2章-练-课时08-T2
\ansitem{7}{A}
\ansnote{详解}{过圆内点的弦以垂直于\(OP\)的弦最短。\(k_{OP}=2\)，故所求斜率\(-1/2\)：\(y-2=-(1/2)(x-1)\)，即\(x+2y-5=0\)，故选：A．}
\end{ansblock}

\liB{变式\textbf{8}}{中档(弦长逆求参)}{}{若圆\(x^{2}+y^{2}-2x+4y+m=0\)截直线\(x+y-3=0\)所得弦长为\(2\)，则实数\(m\)的值为\nobreak (\kongwei)}

\bindopt {\fontsize{10.5pt}{\qpTJgLeadLiti}\selectfont \makebox[\dimexpr0.25\linewidth-0.5em\relax][l]{A．\(-1\)}\makebox[\dimexpr0.25\linewidth-0.5em\relax][l]{B．\(-2\)}\makebox[\dimexpr0.25\linewidth-0.5em\relax][l]{C．\(-4\)}\makebox[\dimexpr0.25\linewidth-0.5em\relax][l]{D．\(-31\)}\par}

\begin{ansblock}[2章-练-课时08-T3]
% ans:2章-练-课时08-T3
\ansitem{8}{C}
\ansnote{详解}{圆心\((1,-2)\)，\(r^{2}=5-m\)。\(d=|1-2-3|/\sqrt{2}=2\sqrt{2}\)。半弦长的平方\(=r^{2}-d^{2}\)，即\(1=5-m-8\)，得\(m=-4\)，故选：C．}
\end{ansblock}

\liB{变式\textbf{9}}{简单(弦长范围求参)}{}{直线\(l\)：\(y=x+m\)与圆\(x^{2}+y^{2}=4\)相交于\(A\)，\(B\)两点，若\(|AB|\geqslant 2\sqrt{3}\)，则实数\(m\)的取值范围为\nobreak (\kongwei)}

\bindopt {\fontsize{10.5pt}{\qpTJgLeadLiti}\selectfont \makebox[\dimexpr0.5\linewidth-1em\relax][l]{A．\([-2,2]\)}\makebox[\dimexpr0.5\linewidth-1em\relax][l]{B．\([-\sqrt{2},\sqrt{2}]\)}\par}

\bindopt {\fontsize{10.5pt}{\qpTJgLeadLiti}\selectfont \makebox[\dimexpr0.5\linewidth-1em\relax][l]{C．\([-1,1]\)}\makebox[\dimexpr0.5\linewidth-1em\relax][l]{D．\([-\sqrt{2}/2,\ \sqrt{2}/2]\)}\par}

\begin{ansblock}[2章-练-课时08-T4]
% ans:2章-练-课时08-T4
\ansitem{9}{B}
\ansnote{详解}{\(|AB|\geqslant 2\sqrt{3}\)，即\(d=\sqrt{r^{2}-(|AB|/2)^{2}}\leqslant\sqrt{4-3}=1\)。又\(d=|m|/\sqrt{2}\)，故\(|m|\leqslant\sqrt{2}\)，故选：B．}
\end{ansblock}

\xiaojie{①识别：弦长公式正用（求弦长）、逆用（由弦长求参）、弦长范围求参、圆内点最短弦，判归本题型.}
\bindp {\kaishu ②操作：几何法三步——定圆心与半径、求弦心距\(d\)、弦长\(=2\sqrt{r^{2}-d^{2}}\)；圆内点最短弦垂直于该点与圆心的连线.}
\bindp {\kaishu ③收束：弦长范围转\(d\)的范围再转参数范围；逆求参时回代表圆条件与相交前提检验.}

% ---------- 探究点三 圆的切线与拓展·收难 ----------
\tjdnr{三}{圆的切线方程}{例\textbf{10}}{中档(截距相等·分类计数)}{与圆\(x^{2}+y^{2}-4x+3=0\)相切，且在\(x\)、\(y\)轴上截距相等的直线共有\nobreak (\kongwei)}

\bindopt {\fontsize{10.5pt}{\qpTJgLeadLiti}\selectfont \makebox[\dimexpr0.25\linewidth-0.5em\relax][l]{A．\(1\)条}\makebox[\dimexpr0.25\linewidth-0.5em\relax][l]{B．\(2\)条}\makebox[\dimexpr0.25\linewidth-0.5em\relax][l]{C．\(3\)条}\makebox[\dimexpr0.25\linewidth-0.5em\relax][l]{D．\(4\)条}\par}

{\ansblockgrayfalse
\begin{ansblock}[2章-练-课时08-T6]
% ans:2章-练-课时08-T6
\ansitem{10}{D}
\ansnote{详解}{圆\((x-2)^{2}+y^{2}=1\)，圆心\((2,0)\)，\(r=1\)。①过原点（两截距同为0）：\(y=kx\)，\(d=|2k|/\sqrt{1+k^{2}}=1\)，即\(k^{2}=1/3\)，2条；②不过原点：\(x/a+y/a=1\)即\(x+y=a\)，\(d=|2-a|/\sqrt{2}=1\)，\(a=2\pm\sqrt{2}\)（均非0），2条。\par\noindent 共4条，故选：D．}
\end{ansblock}}

\liB{变式\textbf{11}}{中档(切两轴定圆)}{}{若过点\((2,1)\)的圆与两坐标轴都相切，则圆心到直线\(2x-y-3=0\)的距离为\nobreak (\kongwei)}

\bindopt {\fontsize{10.5pt}{\qpTJgLeadLiti}\selectfont \makebox[\dimexpr0.5\linewidth-1em\relax][l]{A．\(\sqrt{5}/5\)}\makebox[\dimexpr0.5\linewidth-1em\relax][l]{B．\(2\sqrt{5}/5\)}\par}

\bindopt {\fontsize{10.5pt}{\qpTJgLeadLiti}\selectfont \makebox[\dimexpr0.5\linewidth-1em\relax][l]{C．\(3\sqrt{5}/5\)}\makebox[\dimexpr0.5\linewidth-1em\relax][l]{D．\(4\sqrt{5}/5\)}\par}

{\ansblockgrayfalse
\begin{ansblock}[2章-练-课时08-T7]
% ans:2章-练-课时08-T7
\ansitem{11}{B}
\ansnote{详解}{圆与两轴都切且过\((2,1)\)（第一象限），则圆心\((a,a)\)，\(r=a\)，\(a>0\)。\((2-a)^{2}+(1-a)^{2}=a^{2}\)，即\(a^{2}-6a+5=0\)，\(a=1\)或\(a=5\)。\par\noindent \(a=1\)：圆心\((1,1)\)，\(d=|2-1-3|/\sqrt{5}=2/\sqrt{5}=2\sqrt{5}/5\)；\(a=5\)：圆心\((5,5)\)，\(d=|10-5-3|/\sqrt{5}=2/\sqrt{5}\)，相同。故\(d=2\sqrt{5}/5\)，故选：B．}
\end{ansblock}}

\liB{变式\textbf{12}}{简单(切线长最小值)}{}{过直线\(y=2x-3\)上的点作圆\(C\)：\(x^{2}+y^{2}-4x+6y+12=0\)的切线，则切线长的最小值为\nobreak (\kongwei)}

\bindopt {\fontsize{10.5pt}{\qpTJgLeadLiti}\selectfont \makebox[\dimexpr0.5\linewidth-1em\relax][l]{A．\(\sqrt{55}/5\)}\makebox[\dimexpr0.5\linewidth-1em\relax][l]{B．\(\sqrt{19}\)}\par}

\bindopt {\fontsize{10.5pt}{\qpTJgLeadLiti}\selectfont \makebox[\dimexpr0.5\linewidth-1em\relax][l]{C．\(2\sqrt{5}\)}\makebox[\dimexpr0.5\linewidth-1em\relax][l]{D．\(\sqrt{21}\)}\par}

\begin{ansblock}[2章-练-课时08-T9]
% ans:2章-练-课时08-T9
\ansitem{12}{A}
\ansnote{详解}{圆\((x-2)^{2}+(y+3)^{2}=1\)，圆心\(C(2,-3)\)，\(r=1\)。切线长\(=\sqrt{|PC|^{2}-r^{2}}\)，最小即\(|PC|\)最小\(=C\)到直线距离\(=|4+3-3|/\sqrt{5}=4/\sqrt{5}\)。\par\noindent 最小切线长\(=\sqrt{16/5-1}=\sqrt{11/5}=\sqrt{55}/5\)，故选：A．}
\end{ansblock}

\liB{变式\textbf{13}}{中档(对称轴＋切线长)}{}{已知直线\(l\)：\(x+ay-1=0\)是圆\(C\)：\(x^{2}+y^{2}-6x-2y+1=0\)的对称轴，过点\(A(-1,a)\)作圆\(C\)的一条切线，切点为\(B\)，则\(|AB|=\)\nobreak (\kongwei)}

\bindopt {\fontsize{10.5pt}{\qpTJgLeadLiti}\selectfont \makebox[\dimexpr0.25\linewidth-0.5em\relax][l]{A．\(1\)}\makebox[\dimexpr0.25\linewidth-0.5em\relax][l]{B．\(2\)}\makebox[\dimexpr0.25\linewidth-0.5em\relax][l]{C．\(4\)}\makebox[\dimexpr0.25\linewidth-0.5em\relax][l]{D．\(8\)}\par}

\begin{ansblock}[2章-练-课时08-T10]
% ans:2章-练-课时08-T10
\ansitem{13}{C}
\ansnote{详解}{圆\((x-3)^{2}+(y-1)^{2}=9\)，圆心\((3,1)\)，\(r=3\)。对称轴过圆心：\(3+a-1=0\)，得\(a=-2\)，\(A(-1,-2)\)。\(|AC|=\sqrt{16+9}=5\)，切线长\(|AB|=\sqrt{25-9}=4\)，故选：C．}
\end{ansblock}

\liB{变式\textbf{14}}{中档(切线长)}{}{若斜率为\(\sqrt{3}\)的直线与\(y\)轴交于点\(A\)，与圆\(x^{2}+(y-1)^{2}=1\)相切于点\(B\)，则\(|AB|=\)\kongbai{}．}
\xiexwei{9mm}

\begin{ansblock}[2章-练-课时08-T14]
% ans:2章-练-课时08-T14
\ansitem{14}{√3}
\ansnote{详解}{设直线\(y=\sqrt{3}x+b\)，\(A(0,b)\)。相切：\(d=|b-1|/2=1\)，得\(b=3\)或\(b=-1\)，\(|AC|=2\)（\(C(0,1)\)）。\(|AB|=\sqrt{|AC|^{2}-r^{2}}=\sqrt{4-1}=\sqrt{3}\)．}
\end{ansblock}

\tjdnr{三}{圆的切线方程}{例\textbf{15}}{中档(切线＋弦长综合)}{已知点\(M(3,1)\)，直线\(ax-y+4=0\)及圆\(C\)：\((x-1)^{2}+(y-2)^{2}=4\)．}

\bindp (1) 求过点\(M\)的圆\(C\)的切线方程；

\bindp (2) 若直线\(ax-y+4=0\)与圆\(C\)相切，求实数\(a\)的值；

\bindp (3) 若直线\(ax-y+4=0\)与圆\(C\)相交于\(A\)、\(B\)两点，且弦\(AB\)的长为\(2\sqrt{3}\)，求\(a\)的值.
\xiexwei{18mm}

{\ansblockgrayfalse
\begin{ansblock}[2章-练-课时08-T15]
% ans:2章-练-课时08-T15
\ansitem{15}{(1) x=3 或 3x−4y−5=0；(2) a=0 或 a=4/3；(3) a=−3/4}
\ansnote{详解}{\(C(1,2)\)，\(r=2\)；\(|MC|=\sqrt{4+1}=\sqrt{5}>2\)，\(M\)在圆外。\par\noindent (1) 斜率不存在时\(x=3\)：\(d=|3-1|=2=r\)，是切线。斜率存在时设\(y-1=k(x-3)\)，即\(kx-y-3k+1=0\)：\(d=|k-2-3k+1|/\sqrt{k^{2}+1}=|2k+1|/\sqrt{k^{2}+1}=2\)，即\(4k^{2}+4k+1=4k^{2}+4\)，得\(k=3/4\)，切线\(3x-4y-5=0\)。\par\noindent (2) \(d=|a-2+4|/\sqrt{a^{2}+1}=|a+2|/\sqrt{a^{2}+1}=2\)，即\(a^{2}+4a+4=4a^{2}+4\)，\(3a^{2}-4a=0\)，得\(a=0\)或\(a=4/3\)。\par\noindent (3) 弦长\(2\sqrt{3}\)，半弦\(\sqrt{3}\)，\(d=\sqrt{4-3}=1\)：\(|a+2|/\sqrt{a^{2}+1}=1\)，即\(a^{2}+4a+4=a^{2}+1\)，得\(a=-3/4\)．}
\end{ansblock}}

\liB{变式\textbf{16}}{中档(实际应用·进入时间)}{}{为了保证我国东海油气田海域的海上平台的生产安全，海事部门在某平台\(O\)的正东方向设立了两个观测站\(A\)、\(B\)（点\(A\)在点\(O\)、点\(B\)之间），它们到平台\(O\)的距离分别为\(3\)海里和\(12\)海里，记海平面上到两观测站距离\(PA\)、\(PB\)之比为\(1/2\)的点\(P\)的轨迹为曲线\(E\)，规定曲线\(E\)及其内部区域为安全预警区．}

\bindp (1) 以\(O\)为坐标原点，\(AB\)所在直线为\(x\)轴建立平面直角坐标系，求曲线\(E\)的方程；

\bindp (2) 某日在观测站\(B\)处发现，在该海上平台正南\(2\sqrt{11}\)海里的\(C\)处，有一艘轮船正以每小时\(10\)海里的速度向北偏东\(30^{\circ}\)方向航行，如果航向不变，该轮船是否会进入安全预警区？如果不进入，说明理由；如果进入，则它在安全预警区中的航行时间是几小时．
\xiexwei{18mm}

{\ansblockgrayfalse
\begin{ansblock}[2章-练-课时08-T16]
% ans:2章-练-课时08-T16
\ansitem{16}{(1) x²+y²=36；(2) 一定进入，航行时间 1 小时}
\ansnote{详解}{(1) \(O\)为原点，\(A(3,0)\)，\(B(12,0)\)。由\(|PA|/|PB|=1/2\)：\(4[(x-3)^{2}+y^{2}]=(x-12)^{2}+y^{2}\)，展开：\(4x^{2}-24x+36+4y^{2}=x^{2}-24x+144+y^{2}\)，即\(3x^{2}+3y^{2}=108\)，得\(x^{2}+y^{2}=36\)。\par\noindent (2) \(C(0,-2\sqrt{11})\)。北偏东\(30^{\circ}\)即航线倾角\(60^{\circ}\)，斜率\(\sqrt{3}\)，航线：\(y+2\sqrt{11}=\sqrt{3}x\)，即\(\sqrt{3}x-y-2\sqrt{11}=0\)。\(d=|2\sqrt{11}|/\sqrt{3+1}=\sqrt{11}<6\)，相交，必进入。\par\noindent 弦长\(=2\sqrt{36-11}=10\)，\(t=10/10=1\)（小时）．}
\end{ansblock}}

\xiaojie{①识别：切线方程（过圆上点、过圆外点、切两轴定圆）、切线长与切线条数计数、切线＋弦长综合、阿氏圆实际应用，判归本题型.}
\bindp {\kaishu ②操作：过圆上点切线唯一（垂直于过切点的半径）；过圆外点设斜率必查不存在的支；切线长\(=\sqrt{|PC|^{2}-r^{2}}\)；计数与最值先定圆心半径再分类.}
\bindp {\kaishu ③收束：切线分类不重不漏（过原点/不过原点、斜率存在/不存在）；实际应用建系后全部翻译为距离与弦长，作答带单位与结论语.}

% ============================================================
% 课堂评价 G5（恒5＝3单选+2填空＝导槽 G1~G5；ketang 新制顺排可跨栏，禁 \vbox 装栏）
% 印面号 17—21；多选门：本片正文多选 0 ≤2 合规
% ============================================================
\begin{ketang}{本组共5题（第17—21题），17—19为单选，20—21为填空．}

\jiancestem{{\fontsize{11.4pt}{14pt}\selectfont\heihao {\numboldjian 17}．}\hspace{4.1mm}\tieside{中档(弦与圆心角)}\hspace{2.2mm}圆\(x^{2}+y^{2}=4\)被直线\(y=\sqrt{3}x+2\)截得的劣弧所对的圆心角的大小为\nobreak(\kongwei)}

\bindopt {\fontsize{10.5pt}{\qpTJgLeadPingjia}\selectfont \makebox[\dimexpr0.25\linewidth-0.5em\relax][l]{A．\(30^{\circ}\)}\makebox[\dimexpr0.25\linewidth-0.5em\relax][l]{B．\(60^{\circ}\)}\makebox[\dimexpr0.25\linewidth-0.5em\relax][l]{C．\(90^{\circ}\)}\makebox[\dimexpr0.25\linewidth-0.5em\relax][l]{D．\(120^{\circ}\)}\par}

\begin{ansblock}[2章-导-课时08-G1]
% ans:2章-导-课时08-G1
\ansitem{17}{D}
\ansnote{详解}{圆心\(O(0,0)\)，\(r=2\)。\(O\)到直线距离\(d=|2|/\sqrt{3+1}=1\)。弦心距\(OM=1\)，\(|AM|=\sqrt{r^{2}-d^{2}}=\sqrt{3}\)，故\(\sin\angle AOM=\sqrt{3}/2\)，\(\angle AOM=60^{\circ}\)，圆心角\(\angle AOB=120^{\circ}\)，故选：D．}
\end{ansblock}

\jiancestem{{\fontsize{11.4pt}{14pt}\selectfont\heihao {\numboldjian 18}．}\hspace{4.1mm}\tieside{中档(含参判定)}\hspace{2.2mm}直线\(x\sin\theta+y-1=0\)（\(\theta\in\mathbf{R}\)，\(\theta\neq k\pi+\pi/2\)，\(k\in\mathbf{Z}\)）与圆\(x^{2}+y^{2}=1/2\)的位置关系是\nobreak(\kongwei)}

\bindopt {\fontsize{10.5pt}{\qpTJgLeadPingjia}\selectfont \makebox[\dimexpr0.5\linewidth-1em\relax][l]{A．相切}\makebox[\dimexpr0.5\linewidth-1em\relax][l]{B．相离}\par}

\bindopt {\fontsize{10.5pt}{\qpTJgLeadPingjia}\selectfont \makebox[\dimexpr0.5\linewidth-1em\relax][l]{C．相交}\makebox[\dimexpr0.5\linewidth-1em\relax][l]{D．不确定}\par}

\begin{ansblock}[2章-导-课时08-G2]
% ans:2章-导-课时08-G2
\ansitem{18}{B}
\ansnote{详解}{\(d=|-1|/\sqrt{\sin^{2}\theta+1}\)。由\(\theta\neq k\pi+\pi/2\)知\(\sin\theta\neq\pm 1\)，故\(0\leqslant\sin^{2}\theta<1\)，\(d=1/\sqrt{1+\sin^{2}\theta}\in(1/\sqrt{2},\ 1]\)，恒大于\(r=\sqrt{2}/2=1/\sqrt{2}\)，故相离，故选：B．}
\end{ansblock}

\jiancestem{{\fontsize{11.4pt}{14pt}\selectfont\heihao {\numboldjian 19}．}\hspace{4.1mm}\tieside{中档(切线倾斜角)}\hspace{2.2mm}直线\(l\)过点\((\sqrt{3},\ 3)\)且与圆\(C\)：\(x^{2}+(y-2)^{2}=4\)相切，则直线\(l\)的倾斜角的大小为\nobreak(\kongwei)}

\bindopt {\fontsize{10.5pt}{\qpTJgLeadPingjia}\selectfont \makebox[\dimexpr0.5\linewidth-1em\relax][l]{A．\(30^{\circ}\)或\(120^{\circ}\)}\makebox[\dimexpr0.5\linewidth-1em\relax][l]{B．\(60^{\circ}\)或\(150^{\circ}\)}\par}

\bindopt {\fontsize{10.5pt}{\qpTJgLeadPingjia}\selectfont \makebox[\dimexpr0.5\linewidth-1em\relax][l]{C．\(120^{\circ}\)}\makebox[\dimexpr0.5\linewidth-1em\relax][l]{D．\(30^{\circ}\)}\par}

\begin{ansblock}[2章-导-课时08-G3]
% ans:2章-导-课时08-G3
\ansitem{19}{C}
\ansnote{详解}{\((\sqrt{3})^{2}+(3-2)^{2}=4\)，点在圆上，切线唯一。圆心\((0,2)\)，\(k_{OC}=(3-2)/(\sqrt{3}-0)=\sqrt{3}/3\)，切线斜率\(k=-\sqrt{3}\)，倾斜角\(120^{\circ}\)，故选：C．}
\end{ansblock}

\jiancestem{{\fontsize{11.4pt}{14pt}\selectfont\heihao {\numboldjian 20}．}\hspace{4.1mm}\tieside{中档(求参)}\hspace{2.2mm}若直线\(x-y+1=0\)与圆\((x-a)^{2}+y^{2}=4\)有公共点，则实数\(a\)的取值范围为\kongbai{}.}
\xiexwei{7mm}

\begin{ansblock}[2章-导-课时08-G4]
% ans:2章-导-课时08-G4
\ansitem{20}{[−1−2√2, −1+2√2]}
\ansnote{详解}{圆心\((a,0)\)，\(r=2\)。有公共点\(\iff d=|a+1|/\sqrt{2}\leqslant 2\)，即\(|a+1|\leqslant 2\sqrt{2}\)，得\(-1-2\sqrt{2}\leqslant a\leqslant -1+2\sqrt{2}\)．}
\end{ansblock}

\jiancestem{{\fontsize{11.4pt}{14pt}\selectfont\heihao {\numboldjian 21}．}\hspace{4.1mm}\tieside{中档(中点弦)}\hspace{2.2mm}若点\(P(2,-1)\)为圆\((x-1)^{2}+y^{2}=25\)的弦\(AB\)的中点，则直线\(AB\)的方程是\kongbai{}．}
\xiexwei{7mm}

\begin{ansblock}[2章-导-课时08-G5]
% ans:2章-导-课时08-G5
\ansitem{21}{x−y−3=0}
\ansnote{详解}{圆心\(C(1,0)\)。\(k_{CP}=(-1-0)/(2-1)=-1\)。\(CP\perp AB\)，故\(k_{AB}=1\)。\(y+1=1\cdot(x-2)\)，即\(x-y-3=0\)．}
\end{ansblock}

\end{ketang}

% ---- 尾块（\tailfill 置 \end{multicols} 前；无参＝内容尾随框，件侧不擅自定高） ----
\tailfill

\end{multicols}
\end{document}
